% ===============================================================
% LH-06d: Mi ciudad - Wiederholung + Stundenleistung (LEHRERHINWEISE)
% Spanisch Klasse 7 - Sequenz 6: Mi ciudad (ABSCHLUSS)
% ===============================================================

\documentclass[11pt, a4paper]{article}

\newif\ifswmodus
\swmodusfalse

\usepackage[utf8]{inputenc}
\usepackage[T1]{fontenc}
\usepackage[ngerman]{babel}
\usepackage{helvet}
\renewcommand{\familydefault}{\sfdefault}

\usepackage[margin=2cm]{geometry}
\usepackage{enumitem}
\usepackage{tabularx}
\usepackage{booktabs}
\usepackage{longtable}

\usepackage{xcolor}
\usepackage{amssymb}
\usepackage{tcolorbox}
\tcbuselibrary{skins}

\usepackage{fancyhdr}
\usepackage{lastpage}
\usepackage{needspace}

% FARBEN
\definecolor{spanisch-color}{HTML}{c41e3a}
\definecolor{grau-color}{HTML}{888888}
\definecolor{hellgrau-color}{HTML}{f5f5f5}
\definecolor{basis-color}{HTML}{2a7a4b}
\definecolor{standard-color}{HTML}{2c5aa0}
\definecolor{erweiterung-color}{HTML}{7b1fa2}
\definecolor{loesung-color}{HTML}{c62828}

\definecolor{sw-dark}{HTML}{000000}
\definecolor{sw-gray}{HTML}{666666}
\definecolor{sw-white}{HTML}{ffffff}

\ifswmodus
    \colorlet{spanisch}{sw-dark}
    \colorlet{grau}{sw-gray}
    \colorlet{hellgrau}{sw-white}
    \colorlet{basis}{sw-dark}
    \colorlet{standard}{sw-dark}
    \colorlet{erweiterung}{sw-dark}
    \colorlet{loesung}{sw-dark}
\else
    \colorlet{spanisch}{spanisch-color}
    \colorlet{grau}{grau-color}
    \colorlet{hellgrau}{hellgrau-color}
    \colorlet{basis}{basis-color}
    \colorlet{standard}{standard-color}
    \colorlet{erweiterung}{erweiterung-color}
    \colorlet{loesung}{loesung-color}
\fi

\newcommand{\fachfarbe}{spanisch}

\newcommand{\thema}{Mi ciudad -- Wiederholung + Stundenleistung}
\newcommand{\sequenz}{06}
\newcommand{\block}{d}
\newcommand{\fach}{Spanisch}
\newcommand{\klasse}{7}
\newcommand{\dauer}{90 Min. (Doppelstunde)}
\newcommand{\lerngruppengroesse}{24}
\newcommand{\datum}{\today}

\pagestyle{fancy}
\fancyhf{}
\renewcommand{\headrulewidth}{0.4pt}
\renewcommand{\footrulewidth}{0.4pt}

\lhead{\fach{} -- Klasse \klasse}
\chead{\textbf{LH-\sequenz\block{} (Lehrerhinweise)}}
\rhead{\datum}

\lfoot{}
\cfoot{}
\rfoot{Seite \thepage\ von \pageref{LastPage}}

\newcommand{\B}{\textcolor{basis}{\textbf{[B]}}}
\newcommand{\Std}{\textcolor{standard}{\textbf{[S]}}}
\newcommand{\E}{\textcolor{erweiterung}{\textbf{[E]}}}
\newcommand{\loesung}[1]{\textcolor{loesung}{\textbf{#1}}}
\newcommand{\es}[1]{\textcolor{\fachfarbe}{\textbf{#1}}}

\newtcolorbox{kurzplan}{
    colback=hellgrau,
    colframe=\fachfarbe,
    colbacktitle=white,
    coltitle=\fachfarbe,
    boxrule=1pt,
    arc=0pt,
    fonttitle=\bfseries,
    attach boxed title to top left={yshift=-2mm, xshift=5mm},
    boxed title style={boxrule=0pt, colback=white},
    title=Kurzplan
}

\newtcolorbox{differenzierung}{
    colback=white,
    colframe=grau,
    colbacktitle=white,
    coltitle=grau,
    boxrule=0.5pt,
    arc=0pt,
    fonttitle=\bfseries,
    attach boxed title to top left={yshift=-2mm, xshift=5mm},
    boxed title style={boxrule=0pt, colback=white},
    title=Differenzierung
}

\newtcolorbox{fehlerbox}{
    colback=white,
    colframe=loesung,
    colbacktitle=white,
    coltitle=loesung,
    boxrule=0.5pt,
    arc=0pt,
    fonttitle=\bfseries,
    attach boxed title to top left={yshift=-2mm, xshift=5mm},
    boxed title style={boxrule=0pt, colback=white},
    title=Häufige Fehler
}

\newtcolorbox{materialbox}{
    colback=white,
    colframe=\fachfarbe,
    colbacktitle=white,
    coltitle=\fachfarbe,
    boxrule=0.5pt,
    arc=0pt,
    fonttitle=\bfseries,
    attach boxed title to top left={yshift=-2mm, xshift=5mm},
    boxed title style={boxrule=0pt, colback=white},
    title=Material-Checkliste
}

\begin{document}

\begin{center}
    {\Large\bfseries\textcolor{\fachfarbe}{Lehrerhinweise: \thema}}\\[0.3em]
    {\normalsize LH-\sequenz\block{} | \dauer{} | \textbf{ABSCHLUSS Sequenz 6}}
\end{center}

\vspace{10pt}

\section*{1. Stundenverlauf (Kurzplan)}

\textbf{1. Stunde (45 Min.) -- Schwerpunkt: Wiederholung}

\begin{tabularx}{\textwidth}{|c|l|X|l|}
\hline
\textbf{Zeit} & \textbf{Phase} & \textbf{Inhalt} & \textbf{Material} \\
\hline
0--5 & Einstieg & ``¿Adónde vas hoy?'' Blitzlicht & -- \\
\hline
5--15 & Wiederholung 1 & Orte-Memory: Bild-Wort-Zuordnung & Tafel \\
\hline
15--25 & Wiederholung 2 & ir + estar: Konjugation wiederholen & AB \\
\hline
25--35 & Übung & PA: Dialog Wegbeschreibung üben & AB \\
\hline
35--45 & Sicherung & LP-06d: Quiz zur Selbstüberprüfung & LP \\
\hline
\end{tabularx}

\vspace{1em}

\textbf{2. Stunde (45 Min.) -- Schwerpunkt: Stundenleistung}

\begin{tabularx}{\textwidth}{|c|l|X|l|}
\hline
\textbf{Zeit} & \textbf{Phase} & \textbf{Inhalt} & \textbf{Material} \\
\hline
0--5 & Aktivierung & Letzte Fragen klären, Ruhe herstellen & -- \\
\hline
5--10 & Vorbereitung & SL erklären, Fragen beantworten & SL \\
\hline
10--35 & \textbf{TEST} & Stundenleistung (25 Min.) & SL \\
\hline
35--40 & Abgabe & Einsammeln & -- \\
\hline
40--45 & Abschluss & Ausblick Sequenz 7, Feedback & -- \\
\hline
\end{tabularx}

\section*{2. Differenzierung}

\begin{differenzierung}
\textbf{Legende:} \B{} Basis (BR) | \Std{} Standard (MR) | \E{} Erweiterung (GY)

\textbf{WICHTIG:} Keine Niveaumarkierungen auf Schülermaterial!
\end{differenzierung}

\vspace{1em}

\begin{tabularx}{\textwidth}{|l|X|X|X|}
\hline
& \B{} \textbf{Basis} & \Std{} \textbf{Standard} & \E{} \textbf{Erweiterung} \\
\hline
\textbf{AB-06d} & Aufg. 1-2 Pflicht & Alle Aufgaben & + freier Dialog \\
\hline
\textbf{SL-06d} & Aufg. 1-2 (8P = 4) & Alle (20P) & Bonusaufgabe \\
\hline
\textbf{Zeit SL} & +5 Min. Zugabe & 25 Min. & 25 Min. \\
\hline
\end{tabularx}

\vspace{1em}

\textbf{Konkrete Hinweise:}
\begin{itemize}
    \item \B{} SuS mit Lernschwäche: Vor dem Test Konjugationstabelle nochmal zeigen
    \item \B{} 12/20 Punkte = Note 4 (ausreichend)
    \item \E{} Schnelle: Dürfen Bonusaufgabe bearbeiten (Stadtplan-Beschreibung)
\end{itemize}

\section*{3. Häufige Fehler}

\begin{fehlerbox}
\begin{tabularx}{\textwidth}{|l|X|X|}
\hline
\textbf{Fehler} & \textbf{Beispiel} & \textbf{Reaktion} \\
\hline
ir vs. estar verwechselt & *``Voy en el parque'' & ir = Bewegung ZU, estar = Position AN \\
\hline
a + el nicht kontrahiert & *``Voy a el cine'' & Regel: a + el = al \\
\hline
estar-Konjugation falsch & *``Yo esta en casa'' & estoy, estás, está, estamos \\
\hline
Präposition de vergessen & *``cerca la tienda'' & cerca \textbf{de}, lejos \textbf{de}, enfrente \textbf{de} \\
\hline
Genus bei Orten falsch & *``la parque'' & el parque (mask.), la plaza (fem.) \\
\hline
\end{tabularx}
\end{fehlerbox}

\section*{4. Material-Checkliste}

\begin{materialbox}
\begin{itemize}[label=$\square$]
    \item AB-\sequenz\block.pdf (\lerngruppengroesse{} Kopien)
    \item ML-\sequenz\block.pdf (1$\times$ für Lehrkraft)
    \item \textbf{SL-\sequenz\block.pdf (\lerngruppengroesse{} Kopien) -- TEST!}
    \item SL-\sequenz\block-ML.pdf (1$\times$ für Korrektur)
    \item Lehrwerk: Qué pasa 1, S. 102--119
    \item Tafelmarker / Kreide
    \item Stoppuhr/Timer für Stundenleistung
    \item Stadtplan-Grafik (optional für Tafel)
\end{itemize}
\end{materialbox}

\section*{5. Lösungen AB-\sequenz\block}

\textbf{Merkkasten Orte:}
\begin{itemize}
    \item la tienda = \loesung{das Geschäft}, el supermercado = \loesung{der Supermarkt}
    \item el parque = \loesung{der Park}, la farmacia = \loesung{die Apotheke}
    \item el cine = \loesung{das Kino}, el hospital = \loesung{das Krankenhaus}
    \item la plaza = \loesung{der Platz}, la estación = \loesung{der Bahnhof}
\end{itemize}

\textbf{Merkkasten Verben:}
\begin{itemize}
    \item ir: yo \loesung{voy}, tú \loesung{vas}, él/ella \loesung{va}, nosotros \loesung{vamos}
    \item estar: yo \loesung{estoy}, tú \loesung{estás}, él/ella \loesung{está}, nosotros \loesung{estamos}
\end{itemize}

\textbf{Merkkasten Wegbeschreibung:}
\begin{itemize}
    \item todo recto = \loesung{geradeaus}, cerca de = \loesung{in der Nähe von}
    \item a la derecha = \loesung{nach rechts}, lejos de = \loesung{weit weg von}
    \item a la izquierda = \loesung{nach links}, al lado de = \loesung{neben}
    \item enfrente de = \loesung{gegenüber von}
\end{itemize}

\textbf{Aufgabe 1 (Zuordnung):} (4P)
\begin{itemize}
    \item el parque = der Park, la tienda = das Geschäft
    \item el cine = das Kino, la farmacia = die Apotheke
\end{itemize}

\textbf{Aufgabe 2 (Verben):} (4P)
\begin{itemize}
    \item a) \loesung{voy}, b) \loesung{estás}, c) \loesung{va}, d) \loesung{estamos}
\end{itemize}

\textbf{Aufgabe 3 (Wegbeschreibung):} (6P)
\begin{itemize}
    \item a) \loesung{al lado de}, b) \loesung{todo recto}, \loesung{a la derecha}
    \item c) \loesung{cerca de}, d) \loesung{enfrente de}
\end{itemize}

\textbf{Aufgabe 4 (Dialog):} (6P)
\begin{itemize}
    \item A: \loesung{Perdona, ¿dónde está el cine?} (2P)
    \item B: \loesung{Sigue todo recto y gira a la izquierda.} (2P)
    \item B: \loesung{No, está cerca.} (1P)
    \item B: \loesung{¡De nada!} (1P)
\end{itemize}

\section*{6. Stundenleistung SL-06d -- Übersicht}

\begin{tabular}{|c|l|c|c|}
\hline
\textbf{Aufg.} & \textbf{Inhalt} & \textbf{Punkte} & \textbf{Diff.} \\
\hline
1 & Zuordnung: Orte Spanisch-Deutsch & 4 & alle \\
\hline
2 & Konjugation: ir und estar & 4 & alle \\
\hline
3 & Lückentext: Wegbeschreibung & 6 & alle \\
\hline
4 & Stadtplan: Wegbeschreibung schreiben & 6 & alle \\
\hline
\multicolumn{2}{|r|}{\textbf{Gesamt:}} & \textbf{20} & \\
\hline
\end{tabular}

\vspace{1em}

\textbf{Notenschlüssel (Klasse 7):}

\begin{tabular}{|c|c|c|c|c|c|c|}
\hline
\textbf{Note} & 1 & 2 & 3 & 4 & 5 & 6 \\
\hline
\textbf{ab \%} & 96 & 80 & 60 & 40 & 20 & <20 \\
\hline
\textbf{ab Punkte} & 19 & 16 & 12 & 8 & 4 & <4 \\
\hline
\end{tabular}

\section*{7. Reflexionsnotizen (nach Durchführung)}

\begin{tabularx}{\textwidth}{|l|X|}
\hline
\textbf{SL-Ergebnisse} & Ø: \rule{2cm}{0.4pt} / Beste: \rule{1cm}{0.4pt} / Schlechteste: \rule{1cm}{0.4pt} \\[1.5em]
\hline
\textbf{Problematische Aufgaben} & \\[2em]
\hline
\textbf{Für Sequenz 7} & \\[2em]
\hline
\end{tabularx}

\end{document}
